\documentclass[letterpaper, 12pt]{article}
\input{tex/_preamble}
\begin{document}
\flushleft\includegraphics[width=0.5\textwidth]{img/home/241017_final_logo_mockup.png}

\cosigsection{Legal considerations}{7 May 2026}

Everyone must follow the law, including authors, institutions, publishers, and sleuths.
This guide introduces some rights you can use, as well as some responsibilities you should be aware of.

\textbf{This guide is for informational purposes only and is not legal advice.
If you want advice applicable to your exact situation, consult a lawyer licensed to practice in your \href{https://dictionary.justia.com/jurisdiction}{jurisdiction}.}
Importantly, while laws are typically written in language that looks straightforward,
they use legal \href{https://dictionary.justia.com/term-of-art}{``terms of art''} whose precise meaning is defined elsewhere. Do not assume that you fully understand your legal rights and responsibilities merely because you have read the text of one law. Lessons from one legal jurisdiction will not necessarily be applicable to another one.

\subsection*{Your rights}

When doing post-publication peer review, you may sometimes want to contact authors or institutions.
In the case of public institutions, you may be able to use the institution's \href{https://en.wikipedia.org/wiki/Freedom_of_information_laws_by_country}{``freedom of information'' laws} to your advantage.
For instance, when \href{https://retractionwatch.com/2023/04/24/a-response-to-a-public-records-request-that-raised-more-questions-than-it-answered/}{investigating grant fraud}, Retraction Watch was able to obtain correspondence between a researcher, the staff of the public university this researched worked at, and the staff of a federal agency. \href{https://doi.org/10.1136/bmj.r2028}{Smoliga and Yang} were able to obtain internal documents from the US government showing that an internal review recommended against the approval of a medical device supposed to protect the brain from injuries.
Note that freedom of information laws typically have various exemptions, such as information that could invade someone's privacy without a good reason.

When dealing with private institutions, you cannot use freedom of information laws,
but you can make use of the general rights you have as a citizen, such as the freedom to make factual remarks in public.
Posting your findings on PubPeer can force an institution to react, and is likely to fall under your jurisdiction's right to free speech, so long as you act responsibly. More on your responsibilities below.

Depending on where you work, your employer may provide legal support.
For instance, some universities have policies to provide legal support to employees for activities carried out as part of their work,
possibly including assisting employees with filing lawsuits against third parties.
If you plan to make use of such regulations, make sure you discuss this with your employer first.

You also have general rights such as a right to a fair trial and an attorney if you are accused of crimes for activities related to post-publication peer review, just as for any other activities.
Sleuths, particularly those that comment frequently, sometimes receive legal threats from the researchers who produced the scientific work under scrutiny. These threats are often non-specific and informal, such as vague insinuations that a commenter has engaged in \href{https://dictionary.justia.com/defamation}{defamation}. Threats also often take the form of a formal \href{https://www.csulb.edu/college-of-business/legal-resource-center/article/cease-and-desist-letters-defined-usage-and}{``cease and desist''} letter demanding that the commenter withdraw their statements or face a lawsuit. Cease and desist letters are typically not legally binding documents but, if a lawsuit is actually filed, they can be used as evidence that the defendant was made aware of the plaintiff's claims yet continued their actions.

If you do not want to rely on these rights, or if you do not trust that they will be enforced fairly due to your personal situation or because of where you live, remember that you can post under a pseudonym using services such as PubPeer. In 2016, an appeals court in the US state of Michigan \href{https://retractionwatch.com/2016/12/07/pubpeer-wins-appeal-court-ruling-unmask-commenters/}{ruled} that PubPeer did not have to reveal the \href{https://en.wikipedia.org/wiki/IP_address}{IP address} of pseudonymous commenters. You can even access PubPeer or other such sites using a \href{https://en.wikipedia.org/wiki/VPN_service}{VPN} or even the \href{https://torproject.org/}{Tor network} to make your browsing habits harder to trace.


\subsection*{Your responsibilities}

Some of the best practices for writing an effective PubPeer comment (see COSIG's \href{https://osf.io/sghaq}{entry on the subject}) also make it less likely that the content of a PPPR comment will be regarded as defamatory or otherwise unprotected under a jurisdiction's free speech laws.
Further best practices that make your comments more effective and may reduce your legal exposure include:

\begin{itemize}
    \item Do not accuse someone of a concrete malicious act or crime. \\
    For instance, using language like ``this person committed fraud'' or ``the editors of this journal are an illegal cartel'' is not advisable.
    \item Focus on facts without attributing intent. \\
    For instance, ``the authors faked this figure with Photoshop'' is unacceptable as you cannot know which author did what and for what reason. Instead, write concrete facts such as ``these two figures are the same but rotated''.
    \item Do not disproportionately target specific members of a larger problematic group. \\
    For instance, if you find many authors who benefited from likely citation manipulation in roughly equal proportion, do not focus on one or two of them, even if it's easier to find information about them specifically.
    \item Do not demean people to make a point, especially if the same effect can be achieved with less drama.
    \item Do not deride or make insinuations about groups of people, especially by ethnicity, gender, or other common \href{https://en.wikipedia.org/wiki/Protected_group}{``protected classes''}.
    \item Do not assume that truth or good faith is a defense in your jurisdiction, check what the law actually says, and ask a lawyer if you're unsure.
\end{itemize}

As part of being factual and unbiased, you should have some form of method for what you do. It can be as simple as ``reading COSIG guides and applying them to papers you encounter''.
If you do a lot of post-publication peer review, keep track of what comments you write about what papers, and set some rules for yourself as to what you investigate.

If you release data as part of an investigation, such as a list of suspicious scientific papers, or screenshots of fraudulent conduct you have identified, it is your responsibility to learn about the privacy laws applicable to your jurisdiction.
Avoid naming individuals in public communications about publication integrity issues, unless one or two individuals are centrally involved and you cannot explain what is happening without mentioning them. Generally, it is sufficient to list paper identifiers (such as DOIs) to link readers to relevant articles.
When sharing screenshots (such as screenshots of paper mill advertisements) publicly, it is usually a good idea to redact individuals' names and phone numbers.

\subsection*{Reacting to legal threats}

As discussed above, if you do a lot of post-publication peer review, and especially if you find problems that relate to a specific set of people such as a citation ring, you may receive a threatening letter from an author or their attorney.

In general, if someone claims you violated their rights but cannot name any specific law you broke, it's unlikely that they will get very far, and in fact most of the time such threats do not go to the lawsuit stage because the would-be plaintiff has no concrete basis for their claims. Some jurisdictions even have laws against so-called SLAPP lawsuits, \href{https://en.wikipedia.org/wiki/Strategic_lawsuit_against_public_participation}{``Strategic Lawsuits Against Public Participation''}, to help expedite their dismissal. When it comes to understanding the merit and severity of a legal threat of any form, there is no substitute for consulting a lawyer.

You may be eligible for help from various private sources that seek to help scientific integrity investigations.
For instance, the \href{https://scientificintegrityfund.org/}{Scientific Integrity Fund} provides money for researchers facing concrete legal threats in the US and Canada.

This guide does \emph{not} discuss every theoretical threat you could face. While you could worry about those, and you can always find a lawyer willing to be paid to make a list of them, you should not be any more worried about convoluted interpretations of laws when doing post-publication peer review than when doing any other public activity such as posting on social media.


\subsection*{Examples of legal threats and lawsuits}

In 2023, social scientist Francesca Gino \href{https://www.science.org/content/article/honesty-researcher-facing-fraud-concerns-sues-harvard-and-accusers-25-million}{sued} scientists Uri Simonsohn, Leif Nelson, and Joe Simmons in a US federal court. The three had alleged data manipulation in several papers co-authored by Gino on their blog \href{https://datacolada.org/}{``Data Colada''}. Gino was suspended by her employer, Harvard University, in 2023 and \href{https://www.theguardian.com/education/2025/may/27/harvard-professor-on-leave-falsified-ethics-data}{fired in 2025}. Gino alleged that the three bloggers, Harvard University, and the dean of Harvard Business School (also named as defendant in the lawsuit) had defamed her. Gino sought 25 million USD in damages from these parties. The lawsuit against the bloggers was \href{https://doi.org/10.1126/science.zciy6ft}{dismissed in 2024} with the judge allowing some claims against Harvard to proceed. However, \href{https://www.thecrimson.com/article/2025/7/12/gino-data-colada-legal-fees/}{in 2025, the same judge ruled} that Gino did not need to cover the legal expenses the bloggers had incurred during the proceedings.

In 2017, cancer researcher Carlo Croce sued biologist David Sanders, who had been interviewed by the \textit{New York Times} around apparently falsified data in Croce's articles, for defamation. Croce \href{https://retractionwatch.com/2020/05/13/cancer-researcher-loses-defamation-suit-against-critic/}{lost}, but the judge did not grant Sanders the benefit of Indiana's ``anti-SLAPP'' law either.

In the UK, known to have plaintiff-friendly defamation laws, businessman Arron Banks \href{https://www.bbc.com/news/uk-65644475}{sued} journalist Carole Cadwalladr due to her accusing him of lying about ``his relationship with the Russian government''.
After appealing, Banks' claims were partly upheld, and Cadwalladr has been ordered to pay £1.2m.
As of the time of writing this guide, this case is still under appeal at the level of the European Court of Human Rights.

In 2021, French microbiologist Didier Raoult and French structural biologist Eric Chabriere \href{https://doi.org/10.1038/d41586-021-01430-z}{filed a complaint} with a French public prosecutor accusing US-based sleuth Elisabeth Bik of harassment, extortion and blackmail after Bik left PubPeer comments on dozens of articles authored by Raoult and Chabriere. The public prosecutor \href{https://www.chemistryworld.com/news/integrity-specialist-has-no-case-to-answer-over-blackmail-extortion-allegations-french-officials-find/4019469.article}{informed Bik in 2024} that there was not sufficient evidence to investigate these allegations, effectively ending this legal threat. Raoult made a similar complaint against French teacher Guillaume Limousin. A judge \href{https://www.lexpress.fr/sciences-sante/guillaume-limousin-le-petit-prof-contre-didier-raoult-on-peut-tous-faire-bouger-les-choses-7XXE77HMTNBGVI4D6O5UDBVO7Y/}{dismissed this complaint in 2024} and ordered Raoult to pay for the legal fees Limousin incurred during the proceedings.

In 2024, lawyers representing American chemist Chad Mirkin \href{https://retractionwatch.com/2024/07/16/kavli-prize-winner-threatens-to-sue-critic-for-defamation/}{sent a cease and desist letter} to French physicist Rapha\"{e}l L\'{e}vy alleging defamation. L\'{e}vy had written \href{https://doi.org/10.5281/zenodo.12746733}{a letter} criticizing claims made in \href{https://doi.org/10.1073/pnas.2306973121}{a ``perspectives'' article published by Mirkin}.

\end{document}